\chapter{The formalized Lindeberg principle}

\begin{theorem}[Theorem 1.1, formalized]\label{thm:theorem-one-one}
\lean{Lindeberg.theoremOneOne_paper_statement}
\uses{thm:independent-coupling, lem:product-transfer, lem:error-transfer, lem:third-moment-integrability}
\leanok
Let \((\Omega,\mu)\) be a probability space, let
\(X,Y:\Omega\to(\operatorname{Fin} n\to\mathbb{R})\) have measurable
coordinates, and assume the \(Y_i\)'s are mutually independent.  Suppose
\[
  \int |X_i|^3\,d\mu<\infty,\qquad
  \int |Y_i|^3\,d\mu<\infty,
\]
and
\[
  \int |X_i|^3\,d\mu+\int |Y_i|^3\,d\mu\le M_3
  \quad\text{for every } i.
\]
If \(f:\mathbb{R}^n\to\mathbb{R}\) is \(C^3\) and its \(r\)-fold coordinate
partials satisfy
\[
  |\partial_i^r f(x)|\le L_r
  \qquad (r=1,2,3),
\]
then, with \(A_i\) and \(B_i\) defined using the past sigma-algebra generated
by \(X_j\) for \(j<i\),
\[
\begin{aligned}
  \left|\int_\Omega f(X(\omega))\,d\mu(\omega)
  -\int_\Omega f(Y(\omega))\,d\mu(\omega)\right|
  &\le
  \sum_{i:\operatorname{Fin} n}
    \left(A_iL_1+\frac12B_iL_2\right)
  +\frac16 nL_3M_3 .
\end{aligned}
\]
No independence between \(X\) and \(Y\) is assumed in this public statement.
\end{theorem}

This blueprint tracks how \texttt{Lindeberg/Basic.lean} proves the theorem
above.  The comparison vector \(Y\) has independent coordinates, while Lean
removes any auxiliary independence assumption between \(X\) and \(Y\) by
passing to the product probability space and then transferring the moment
errors back to the original space.

The Lean development uses zero-based coordinates \(i : \operatorname{Fin} n\).
Thus the paper's sigma-algebra generated by \(X_1,\ldots,X_{i-1}\) becomes
the sigma-algebra generated by the coordinates \(j<i\).

\section{Deterministic finite-sum bounds}

\begin{definition}[Coordinate and global bounds]\label{def:deterministic-bounds}
\lean{Lindeberg.coordinateBound, Lindeberg.theoremOneOneBound}
\leanok
For functions \(A,B : \operatorname{Fin} n \to \mathbb{R}\) and constants
\(L_1,L_2,L_3,M_3\), Lean defines
\[
  c_i=A_iL_1+\frac12 B_iL_2+\frac16 L_3M_3
\]
as \texttt{coordinateBound}, and the displayed right hand side of Theorem 1.1
as
\[
  \sum_i \left(A_iL_1+\frac12 B_iL_2\right)
  + \frac16 nL_3M_3 .
\]
\end{definition}

\begin{lemma}[Sum of coordinate bounds]\label{lem:sum-coordinate-bound}
\lean{Lindeberg.sum_coordinateBound_eq_theoremOneOneBound}
\uses{def:deterministic-bounds}
\leanok
The finite sum of the per-coordinate bounds is exactly
\texttt{theoremOneOneBound}.
\end{lemma}

\begin{theorem}[Deterministic telescope]\label{thm:deterministic-telescope}
\lean{Lindeberg.theoremOneOne_deterministic, Lindeberg.theoremOneOne_deterministic_of_thirdMomentBound}
\uses{def:deterministic-bounds, lem:sum-coordinate-bound}
\leanok
Suppose real numbers \(e_X,e_Y\) and coordinate increments \(\Delta_i\)
satisfy
\[
  e_X-e_Y=\sum_i \Delta_i
\]
and each \(|\Delta_i|\) is bounded by the coordinate Taylor/Lindeberg error.
Then
\[
  |e_X-e_Y|\le \texttt{theoremOneOneBound}(A,B,L_1,L_2,L_3,M_3).
\]
Lean also includes the variant where the third-moment term depends on \(i\)
and is bounded coordinatewise by \(M_3\).
\end{theorem}

\section{Probabilistic quantities}

\begin{definition}[Conditional moment errors]\label{def:moment-errors}
\lean{Lindeberg.expectation, Lindeberg.firstMomentError, Lindeberg.secondMomentError, Lindeberg.thirdMomentContribution}
\leanok
For a measure \(\mu\), a family of sub-sigma-algebras \(\mathcal F_i\), and
coordinate processes \(X_i,Y_i\), the formalized errors are
\[
  A_i =
  \int_\Omega
    \left|
      \operatorname{condExp}_{\mathcal F_i}(X_i)(\omega)
      - \int_\Omega Y_i\,d\mu
    \right|\,d\mu(\omega),
\]
\[
  B_i =
  \int_\Omega
    \left|
      \operatorname{condExp}_{\mathcal F_i}(X_i^2)(\omega)
      - \int_\Omega Y_i^2\,d\mu
    \right|\,d\mu(\omega),
\]
and the coordinate third-moment contribution is
\[
  \int_\Omega |X_i|^3\,d\mu+\int_\Omega |Y_i|^3\,d\mu .
\]
\end{definition}

\begin{lemma}[Third moments give lower moments]\label{lem:third-moment-integrability}
\lean{Lindeberg.abs_le_abs_cube_add_one, Lindeberg.abs_sq_le_abs_cube_add_one, Lindeberg.integrable_of_integrable_abs_cube, Lindeberg.integrable_sq_of_integrable_abs_cube}
\leanok
On a finite measure space, integrability of \(|x|^3\) implies integrability
of \(x\) and \(x^2\).  The proof uses the elementary bounds
\[
  |a|\le |a|^3+1,
  \qquad
  |a|^2\le |a|^3+1 .
\]
\end{lemma}

\section{The Lindeberg interpolation}

\begin{definition}[Coordinate processes and sigma-algebras]\label{def:sigma-algebras}
\lean{Lindeberg.coordinateProcess, Lindeberg.vectorSigma, Lindeberg.pastSigma, Lindeberg.futureSigma, Lindeberg.frozenSigma}
\leanok
Lean separates a random vector into coordinate random variables, defines the
full coordinate sigma-algebra \(\sigma(X_0,\ldots,X_{n-1})\), and defines
\[
  \mathcal P_i=\sigma(X_j:j<i),\qquad
  \mathcal U_i=\sigma(Y_j:i<j).
\]
The frozen sigma-algebra is
\[
  \mathcal G_i=\mathcal P_i\vee\mathcal U_i .
\]
\end{definition}

\begin{definition}[Interpolated and frozen vectors]\label{def:interpolation}
\lean{Lindeberg.lindebergInterpolate, Lindeberg.lindebergFrozen}
\uses{def:sigma-algebras}
\leanok
The interpolation vector is
\[
  Z_k(\omega)_j=
  \begin{cases}
    X(\omega)_j, & j<k,\\
    Y(\omega)_j, & k\le j,
  \end{cases}
\]
with \(Z_0=Y\) and \(Z_n=X\).  The frozen vector for coordinate \(i\) is
\[
  Z_i^0(\omega)_j=
  \begin{cases}
    X(\omega)_j, & j<i,\\
    0, & j=i,\\
    Y(\omega)_j, & i<j.
  \end{cases}
\]
\end{definition}

\begin{lemma}[Measurability of the interpolation data]\label{lem:interpolation-measurable}
\lean{Lindeberg.pastSigma_le, Lindeberg.futureSigma_le, Lindeberg.vectorSigma_le, Lindeberg.pastSigma_le_vectorSigma, Lindeberg.futureSigma_le_vectorSigma, Lindeberg.coordinate_stronglyMeasurable_vectorSigma, Lindeberg.frozenSigma_le, Lindeberg.pastSigma_le_frozenSigma, Lindeberg.lindebergInterpolate_measurable, Lindeberg.lindebergFrozen_measurable_frozenSigma}
\uses{def:sigma-algebras, def:interpolation}
\leanok
The coordinate and frozen sigma-algebras are sub-sigma-algebras of the ambient
measurable space whenever the coordinate random variables are measurable.  The
frozen vector \(Z_i^0\) is measurable with respect to \(\mathcal G_i\).
\end{lemma}

\begin{lemma}[Replacement by coordinate updates]\label{lem:update-identities}
\lean{Lindeberg.lindebergInterpolate_zero, Lindeberg.lindebergInterpolate_all, Lindeberg.lindebergFrozen_self, Lindeberg.lindebergInterpolate_eq_update_frozen_self, Lindeberg.lindebergInterpolate_succ_eq_update_frozen}
\uses{def:interpolation}
\leanok
The adjacent interpolation vectors are obtained from the same frozen vector by
updating coordinate \(i\): one update inserts \(Y_i\) and the other inserts
\(X_i\).  This is the formal bridge from the vector interpolation to a
one-dimensional Taylor expansion.
\end{lemma}

\begin{definition}[Coordinate slices and derivatives]\label{def:coordinate-derivatives}
\lean{Lindeberg.coordinateSlice, Lindeberg.coordinateTaylorRemainder, Lindeberg.coordinateDerivative, Lindeberg.coordinatePartialDerivative, Lindeberg.frozenCoordinateDerivative}
\uses{def:interpolation}
\leanok
For \(f : \mathbb{R}^n\to\mathbb{R}\), a vector \(z\), and a coordinate \(i\),
the coordinate slice is
\[
  t\mapsto f(z[i:=t]).
\]
Lean defines the corresponding first, second, and third ordinary derivatives,
the paper's coordinate partial derivatives, and the quadratic Taylor
remainder around \(t=0\).
\end{definition}

\begin{lemma}[Coordinate derivative transport]\label{lem:coordinate-derivative-transport}
\lean{Lindeberg.ContDiff.coordinateSlice, Lindeberg.coordinateSlice_update_self, Lindeberg.coordinateDerivative_eq_coordinatePartialDerivative_update_zero, Lindeberg.iteratedDeriv_coordinateSlice_eq_coordinatePartialDerivative_update, Lindeberg.coordinateDerivative_one_eq_lineDeriv_of_coord_zero, Lindeberg.deriv_coordinateSlice_eq_lineDeriv_update, Lindeberg.coordinateDerivative_two_eq_lineDeriv_lineDeriv_of_coord_zero, Lindeberg.frozenCoordinateDerivative_one_stronglyMeasurable, Lindeberg.frozenCoordinateDerivative_two_stronglyMeasurable}
\uses{def:coordinate-derivatives, lem:interpolation-measurable}
\leanok
A \(C^k\) function has \(C^k\) coordinate slices.  The derivatives of those
slices agree with the coordinate partial derivatives used in the paper, and
the first two frozen coordinate derivatives are strongly measurable with
respect to the frozen sigma-algebra.
\end{lemma}

\begin{proposition}[The telescoping identity]\label{prop:lindeberg-telescope}
\lean{Lindeberg.lindeberg_step_eq_coordinateTaylorRemainders, Lindeberg.lindeberg_step_ae_eq_coordinateTaylorRemainders, Lindeberg.lindeberg_integral_telescope}
\uses{def:interpolation, def:coordinate-derivatives, lem:update-identities}
\leanok
The expectation difference telescopes:
\[
  \int f(X)\,d\mu-\int f(Y)\,d\mu
  =
  \sum_i \int \bigl(f(Z_{i+1})-f(Z_i)\bigr)\,d\mu .
\]
For each coordinate, the step is rewritten as the difference of the two
quadratic Taylor expansions through the frozen vector \(Z_i^0\), plus the two
coordinate Taylor remainders.
\end{proposition}

\section{Conditional expectation and independence}

\begin{lemma}[Bounding integrals by conditional moment errors]\label{lem:integral-bound}
\lean{Lindeberg.abs_integral_mul_le_integral_abs_mul_bound}
\uses{def:moment-errors}
\leanok
If a derivative factor \(D\) is almost everywhere bounded by \(L\), then
\[
  \left|\int a(\omega)D(\omega)\,d\mu(\omega)\right|
  \le L\int |a(\omega)|\,d\mu(\omega).
\]
This converts the conditional-expectation identities into the \(A_iL_1\) and
\(B_iL_2\) terms.
\end{lemma}

\begin{lemma}[Pulling out conditional expectations]\label{lem:condexp-pullout}
\lean{Lindeberg.integral_mul_eq_integral_condExp_mul, Lindeberg.integral_sub_mul_eq_condExp_sub_integral_mul, Lindeberg.integral_mul_eq_integral_ae_condExp_mul, Lindeberg.integral_sub_mul_eq_ae_condExp_sub_integral_mul}
\uses{def:moment-errors}
\leanok
For functions measurable with respect to a sub-sigma-algebra, integration
against \(x\) can be replaced by integration against
\(\operatorname{condExp}(x)\).  Lean proves both exact and almost-everywhere
versions, and then applies them to the linear and quadratic terms.
\end{lemma}

\begin{lemma}[Independence of frozen derivatives from future coordinates]\label{lem:frozen-independence}
\lean{Lindeberg.indepFun_of_indep_comap_of_stronglyMeasurable_right, Lindeberg.indepFun_sq_of_indep_comap_of_stronglyMeasurable_right, Lindeberg.y_frozen_derivative_indepFun_of_indep_frozen, Lindeberg.measurableRectangles, Lindeberg.isPiSystem_measurableRectangles, Lindeberg.generateFrom_measurableRectangles_eq_sup, Lindeberg.y_coord_indep_futureSigma_of_iIndep, Lindeberg.indep_left_sup_of_indep_left_right, Lindeberg.y_coord_indep_frozenSigma_of_indep_vectors_iIndep, Lindeberg.indepFun_of_indep_of_stronglyMeasurable, Lindeberg.integral_inter_eq_measure_mul_integral_of_indep}
\uses{def:sigma-algebras, def:coordinate-derivatives}
\leanok
The \(Y_i\) coordinate and its square are independent of the frozen derivative
data when the \(Y\)-coordinates are mutually independent and the \(X\)- and
\(Y\)-generated sigma-algebras are independent.  The proof builds this from
independence of the relevant sigma-algebras and a pi-system description of
rectangles in a supremum sigma-algebra.
\end{lemma}

\begin{lemma}[Reducing frozen conditional expectations to the past]\label{lem:frozen-condexp}
\lean{Lindeberg.condExp_sup_indep_eq_condExp_left, Lindeberg.condExp_frozenSigma_eq_pastSigma_of_indep_vectors, Lindeberg.condExp_frozenSigma_sq_eq_pastSigma_of_indep_vectors, Lindeberg.condExp_eq_of_indep_sub_integral_zero, Lindeberg.integral_sub_condExp_eq_zero, Lindeberg.condExp_eq_condExp_of_indep_sub_condExp}
\uses{def:sigma-algebras, lem:frozen-independence}
\leanok
Although the Taylor derivative is measurable with respect to the larger frozen
sigma-algebra \(\mathcal G_i\), the conditional expectation of \(X_i\) and
\(X_i^2\) against \(\mathcal G_i\) agrees almost everywhere with the
conditional expectation against the past sigma-algebra \(\mathcal P_i\).
\end{lemma}

\section{Taylor estimates}

\begin{lemma}[One-dimensional Taylor remainder]\label{lem:one-dimensional-taylor}
\lean{Lindeberg.taylorWithinEval_two_zero, Lindeberg.coordinateTaylorRemainder_eq_taylorWithin, Lindeberg.oneDim_taylorWithin_remainder_bound_of_pos, Lindeberg.coordinateTaylorRemainder_bound_of_pos, Lindeberg.iteratedDerivWithin_one_Icc_zero_eq_iteratedDeriv, Lindeberg.iteratedDerivWithin_two_Icc_zero_eq_iteratedDeriv, Lindeberg.iteratedDerivWithin_three_Icc_eq_iteratedDeriv_of_mem_Ioo, Lindeberg.oneDim_taylor_remainder_bound_of_pos, Lindeberg.oneDim_taylor_remainder_bound}
\uses{def:coordinate-derivatives}
\leanok
For a \(C^3\) real function \(g\) with \(|g'''|\le L_3\), the quadratic Taylor
remainder around \(0\) satisfies
\[
  \left|g(x)-g(0)-xg'(0)-\frac12x^2g''(0)\right|
  \le \frac16 L_3 |x|^3 .
\]
Lean proves the positive-orientation statement first and then obtains the
two-sided version by applying it to \(u\mapsto g(-u)\).
\end{lemma}

\begin{lemma}[Coordinate Taylor remainder]\label{lem:coordinate-taylor}
\lean{Lindeberg.coordinateTaylorRemainder_bound, Lindeberg.coordinateTaylorRemainder_integrable_of_bound}
\uses{def:coordinate-derivatives, lem:coordinate-derivative-transport, lem:one-dimensional-taylor}
\leanok
Applying the one-dimensional estimate to coordinate slices gives
\[
  |R_i(z,x)|
  \le \frac16 L_3 |x|^3 .
\]
The same bound also supplies integrability of the coordinate Taylor remainder
under the third absolute moment hypotheses.
\end{lemma}

\begin{lemma}[Integrability of the interpolated function]\label{lem:interpolated-integrability}
\lean{Lindeberg.coordinateSlice_abs_sub_zero_le, Lindeberg.abs_apply_le_abs_zero_add_sum, Lindeberg.integrable_f_apply_of_integrable_abs_coords}
\uses{def:coordinate-derivatives, lem:third-moment-integrability}
\leanok
The derivative bound on \(f\) gives a linear growth estimate
\[
  |f(z)|\le |f(0)|+\sum_i L_1|z_i|.
\]
Consequently \(f(Z)\) is integrable whenever the coordinates have integrable
absolute values.
\end{lemma}

\begin{theorem}[One-step Taylor estimate]\label{thm:one-step-estimate}
\lean{Lindeberg.oneStep_taylor_remainder_bound, Lindeberg.oneStep_taylor_condExp_estimate, Lindeberg.integral_taylor_step_expansion, Lindeberg.oneStep_taylor_condExp_estimate_of_independence, Lindeberg.oneStep_taylor_condExp_estimate_of_enlarged_conditioning, Lindeberg.oneStep_taylor_condExp_estimate_complete, Lindeberg.oneStep_taylor_condExp_estimate_complete_enlarged}
\uses{def:moment-errors, lem:integral-bound, lem:condexp-pullout, lem:frozen-independence, lem:frozen-condexp, lem:coordinate-taylor}
\leanok
For a single coordinate replacement, the Taylor expansion, the conditional
expectation identities, and the remainder estimate give
\[
\begin{aligned}
  |\Delta_i|
  &\le
  \left|\int
    \bigl(\operatorname{condExp}_{\mathcal F_i}(X_i)-\mathbb E Y_i\bigr)
    D_{1,i}\,d\mu\right| \\
  &\quad+\frac12
  \left|\int
    \bigl(\operatorname{condExp}_{\mathcal F_i}(X_i^2)-\mathbb E Y_i^2\bigr)
    D_{2,i}\,d\mu\right| \\
  &\quad+\frac16L_3
  \left(\mathbb E|X_i|^3+\mathbb E|Y_i|^3\right).
\end{aligned}
\]
The later variants instantiate the abstract step with the actual frozen
Lindeberg data.
\end{theorem}

\section{The main proof ladder}

\begin{theorem}[Probabilistic summation]\label{thm:probabilistic-summation}
\lean{Lindeberg.theoremOneOne_probability}
\uses{thm:deterministic-telescope, def:moment-errors, lem:integral-bound}
\leanok
If each Lindeberg increment has the conditional-expectation form supplied by
the one-step estimate, then the full expectation difference is bounded by
\[
  \texttt{theoremOneOneBound}(A,B,L_1,L_2,L_3,M_3).
\]
\end{theorem}

\begin{theorem}[Complete abstract probabilistic theorem]\label{thm:probability-complete}
\lean{Lindeberg.theoremOneOne_probability_complete}
\uses{thm:probabilistic-summation, thm:one-step-estimate}
\leanok
The abstract probabilistic theorem packages all per-coordinate hypotheses:
Taylor expansions, integrability, measurability, conditional-expectation
identities, independence, and the third-moment bound.  Its conclusion is the
same \texttt{theoremOneOneBound}.
\end{theorem}

\begin{theorem}[Interpolation theorem]\label{thm:interpolation-complete}
\lean{Lindeberg.theoremOneOne_probability_complete_interpolation, Lindeberg.theoremOneOne_probability_complete_interpolation_residual, Lindeberg.theoremOneOne_probability_complete_interpolation_residual_enlarged}
\uses{prop:lindeberg-telescope, thm:probability-complete, lem:interpolated-integrability, lem:coordinate-taylor}
\leanok
The complete theorem is specialized to the actual Lindeberg interpolation
\(Z_i\) and frozen vectors \(Z_i^0\).  The residual versions insert the
coordinate Taylor remainders explicitly and verify their integrability and
third-moment bounds.
\end{theorem}

\begin{theorem}[Independence specialized to the paper]\label{thm:paper-independence}
\lean{Lindeberg.theoremOneOne_probability_complete_interpolation_residual_of_indep_residuals, Lindeberg.theoremOneOne_probability_complete_interpolation_residual_past_frozen, Lindeberg.theoremOneOne_probability_complete_interpolation_residual_past_frozen_of_indep_vectors, Lindeberg.theoremOneOne_probability_complete_interpolation_residual_past_frozen_of_paper_independence, Lindeberg.theoremOneOne_probability_complete_interpolation_exact_taylor_of_paper_independence}
\uses{thm:interpolation-complete, lem:frozen-independence, lem:frozen-condexp}
\leanok
The proof ladder progressively removes abstract independence hypotheses and
replaces them with the hypotheses appearing in Theorem 1.1: independence of the
\(Y_i\)'s and, for the internal engine only, independence of the sigma-algebras
generated by \(X\) and \(Y\).  The final theorem in this group uses the exact
coordinate Taylor remainders and the paper's derivative bounds.
\end{theorem}

\section{Product coupling and the public theorem}

\begin{lemma}[Product-space transfer facts]\label{lem:product-transfer}
\lean{Lindeberg.iIndep_comap_measurePreserving, Lindeberg.condExp_comap_measurePreserving_ae_eq, Lindeberg.vectorSigma_prod_fst_le, Lindeberg.vectorSigma_prod_snd_le, Lindeberg.pastSigma_prod_fst, Lindeberg.indep_fst_snd_prod, Lindeberg.measurePreserving_fst_prod, Lindeberg.measurePreserving_snd_prod}
\uses{def:sigma-algebras}
\leanok
On the product probability space \(\Omega\times\Omega\), the first and second
projections are measure-preserving, their pulled-back sigma-algebras are
independent, and mutual independence of the \(Y_i\)'s pulls back through the
second projection.  Conditional expectations against a pulled-back
sub-sigma-algebra agree almost everywhere with the original conditional
expectation composed with the projection.
\end{lemma}

\begin{lemma}[Moment errors survive product coupling]\label{lem:error-transfer}
\lean{Lindeberg.firstMomentError_prod_eq, Lindeberg.secondMomentError_prod_eq}
\uses{def:moment-errors, def:sigma-algebras, lem:product-transfer, lem:third-moment-integrability}
\leanok
Let \(X'(p)=X(p_1)\) and \(Y'(p)=Y(p_2)\) on \(\Omega\times\Omega\).  The
first and second conditional moment errors computed for \(X',Y'\) on the
product space are exactly the original \(A_i\) and \(B_i\).
\end{lemma}

\begin{theorem}[Independent-coupling engine]\label{thm:independent-coupling}
\lean{Lindeberg.theoremOneOne_paper_statement_independent_coupling}
\uses{thm:paper-independence, def:moment-errors}
\leanok
Assume \(X\) and \(Y\) are already coupled so that the sigma-algebra generated
by all \(X\)-coordinates is independent of the sigma-algebra generated by all
\(Y\)-coordinates, and assume the \(Y_i\)'s are mutually independent.  Under
the finite third-moment and derivative-bound hypotheses, Lean proves the
paper's displayed estimate:
\[
\begin{aligned}
  \left|\int f(X)\,d\mu-\int f(Y)\,d\mu\right|
  &\le
  \sum_i\left(A_iL_1+\frac12B_iL_2\right)
  +\frac16 nL_3M_3 .
\end{aligned}
\]
\end{theorem}
The public theorem stated at the beginning is obtained by applying this
engine on \(\Omega\times\Omega\), using \(X\) on the first coordinate and
\(Y\) on the second, and then rewriting the product-space errors back to the
original probability space.
